
During my bachelor's degree, not much happened. I was just someone who
followed the book ---what everyone imagines a chemistry student to be,
nothing more---. Ok, one hobby (music), \textit{et pas plus}. not a video-game friki,
not a language nerd, not a chess player, and definitely not a popular
kid.

Academically, there were a few turning points. Maybe the first was
electromagnetism, i failed it, during kinematics and dynamics i was just
solving problems without really understanding them ---just another mandatory
course---. Electromagnetism was different. when i had to properly study to
pass the exam, i don't know how, but i ended up getting 10/10, and i
started to like physics. The second turning point was quantum mechanics,
again, i initially thought ``another mandatory course'', but i loved it. from
that moment on, i knew i wanted to be a theoretical chemist.

During the last part of my bachelor's degree (five years in mexico), i
started taking astrophysics classes. I still love astro-things
(astrocositas). By then, i knew i wanted to do a master's degree, so i
applied to astrophysics in Australia and to theoretical chemistry in Spain
(\href{https://www.emtccm.org}{TCCM} ---an amazing master if you like theoretical chemistry---). I wasn't sure
what to choose, but the decision wasn't entirely mine ---it was in the
hands? of covid---, the pandemic decided for me: i moved to Spain.

A completely new chapter started: postgraduate life. Maybe not the best
financial decision, but definitely a fun one. During the second year of my
master's, i did an erasmus in Paris, it was a really nice time, and i
decided i wanted to stay longer in France ---not Paris, but Rouen---. I decided
to do a PhD, because\ldots{} why not? phd time. And apparently not satisfied with
avoiding a ``real'' job, i decided to do a postdoc.

